\documentclass[../../main.tex]{subfiles}

\begin{document}

\chapter{Array, String, dan Pencetakan Bilangan}
\label{chap:array-string}

% =========================================================
% OUTCOME PEMBELAJARAN KHUSUS BAB
% =========================================================

\begin{subcpmkbox}
\vspace{1.1em}
\SubCPMK{3.6} --- Mengolah array dan string serta mencetak bilangan
desimal dan heksadesimal ke layar.

\vspace{0.4em}
\textbf{CPMK yang didukung:} \CPMK{03} \quad\textbullet\quad
\textbf{CPL yang didukung:} \CPL{03}
\end{subcpmkbox}

\begin{tujuanbox}
\vspace{1.1em}
Setelah mempelajari bab ini, mahasiswa mampu:
\begin{enumerate}
    \item mendeklarasikan array satu dimensi dan multi dimensi;
    \item menghitung offset elemen array;
    \item mengakses array dengan \eng{loop} dan register indeks;
    \item menjelaskan representasi string sebagai array byte;
    \item menyalin, membandingkan, dan mencari isi string;
    \item mengubah bilangan menjadi deret digit desimal;
    \item mengubah bilangan menjadi digit heksadesimal;
    \item mencetak hasil konversi dengan layanan DOS.
\end{enumerate}
\end{tujuanbox}

\begin{materibox}
\vspace{1.1em}
\begin{enumerate}
    \item Array satu dimensi dan penempatan di memori.
    \item Array multi dimensi dan perhitungan offset.
    \item String sebagai array byte dan penanda akhir.
    \item Operasi string: salin, bandingkan, cari.
    \item Konversi bilangan ke digit desimal.
    \item Konversi bilangan ke digit heksadesimal.
    \item Pencetakan karakter dengan layanan DOS.
\end{enumerate}
\end{materibox}

\subfile{section-01}
\subfile{section-02}
\subfile{section-03}

\section{Contoh Program}
\label{sec:contoh-bab15}

Bagian ini menyajikan dua program yang \emph{sengaja dibuat salah}
beserta versi perbaikannya. Keduanya dirakit tanpa kesalahan sintaks,
sehingga assembler tidak memberi petunjuk apa pun: kesalahan baru
tampak ketika program dijalankan dan ditelusuri.

\subsection{Contoh 1: Perulangan Tak Berujung}

Program pada Kode~\ref{lst:bug-loop} bermaksud menjumlahkan bilangan
1 sampai 5, yang seharusnya menghasilkan \hex{000F} atau \dec{15}.
Kenyataannya program tidak pernah selesai, dan karena baris pencetakan
berada setelah perulangan, program tampak ``menggantung'' tanpa
penjelasan apa pun bagi pemakainya.

\lstinputlisting[
    language=MASM,
    caption={Program ber-bug \file{bug-loop.asm}: perulangan tak
    berujung.},
    label={lst:bug-loop}
]{\CodeDir/bug-loop.asm}

Versi perbaikan pada Kode~\ref{lst:bug-loop-fix} menghapus
\asm{INC CX} dan menambahkan rutin pencetakan dua digit desimal agar
hasilnya dapat dilihat di layar.

\lstinputlisting[
    language=MASM,
    caption={Versi perbaikan \file{bug-loop-fix.asm}.},
    label={lst:bug-loop-fix}
]{\CodeDir/bug-loop-fix.asm}

\begin{session}[caption={Keluaran program
\file{bug-loop-fix.exe}.}, label={lst:keluaran-loop-fix}]
Jumlah 1..5 = 15
\end{session}

\begin{infobox}
Perhatikan bahwa \asm{LOOP} sudah melakukan dua pekerjaan sekaligus:
mengurangi \reg{CX} sebesar satu, lalu melompat bila \reg{CX} belum
nol. Karena itu, penghitung perulangan \emph{tidak perlu} dan
\emph{tidak boleh} diubah sendiri di dalam badan perulangan.
\end{infobox}

\subsection{Contoh 2: Stack yang Tidak Seimbang}

Program pada Kode~\ref{lst:bug-stack} memanggil prosedur
\asm{HITUNG}, lalu bermaksud mencetak pesan penutup. Program berhasil
dirakit, tetapi ketika dijalankan kendali melompat ke alamat yang
tidak dikenal dan pesan penutup tidak pernah muncul. Prosedur
tersebut menekan tiga register (\reg{AX}, \reg{BX}, \reg{CX}) tetapi
hanya memulihkan dua, sehingga \asm{RET} tidak menemukan alamat
kembali pada puncak \eng{stack}.

\lstinputlisting[
    language=MASM,
    caption={Program ber-bug \file{bug-stack.asm}: \eng{stack} tidak
    seimbang.},
    label={lst:bug-stack}
]{\CodeDir/bug-stack.asm}

Versi perbaikan pada Kode~\ref{lst:bug-stack-fix} menambahkan satu
instruksi \asm{POP AX} sehingga jumlah \asm{PUSH} dan \asm{POP}
kembali seimbang.

\lstinputlisting[
    language=MASM,
    caption={Versi perbaikan \file{bug-stack-fix.asm}.},
    label={lst:bug-stack-fix}
]{\CodeDir/bug-stack-fix.asm}

\begin{session}[caption={Keluaran program
\file{bug-stack-fix.exe}.}, label={lst:keluaran-stack-fix}]
Selesai.
\end{session}

\begin{peringatanbox}
Kaidah yang harus dihafal: \emph{setiap} \asm{PUSH} memerlukan
\asm{POP} pasangannya pada jalur eksekusi yang sama. Pada prosedur
yang memakai beberapa jalur keluar, pastikan seluruh jalur
menyeimbangkan \eng{stack} sebelum \asm{RET}. Bila tidak, \asm{RET}
akan melompat ke alamat yang salah dan gejalanya muncul jauh dari
penyebabnya.
\end{peringatanbox}

\section{Praktikum}
\label{sec:praktikum-bab15}

\begin{praktikumbox}
\textbf{Judul:} \eng{Debugging} terpandu tiga kasus bug tingkat
instruksi.\\[0.3em]
\textbf{Sub-CPMK:} 4.1 dan 4.2 \quad\textbullet\quad
\textbf{Alokasi waktu:} 100 menit \quad\textbullet\quad
\textbf{Bentuk:} individu, laporan tertulis
\end{praktikumbox}

\subsection{Tujuan Praktikum}

Setelah menyelesaikan praktikum ini, mahasiswa mampu menelusuri
program pada level instruksi dengan perintah \code{R}, \code{U},
\code{T}, \code{P}, \code{G}, dan \code{D}; mengisi tabel \eng{trace}
register dan \eng{flag}; membaca \eng{dump} memori dan \eng{stack};
serta menemukan dan memperbaiki perulangan tak berujung, kesalahan
\eng{off-by-one}, dan \eng{stack} yang tidak seimbang.

\subsection{Alat dan Bahan}

\begin{itemize}
    \item Komputer dengan sistem operasi Windows.
    \item DOSBox beserta berkas \file{DEBUG.EXE}.
    \item TASM dan TLINK (atau Emu8086 sebagai pembanding).
    \item Berkas \file{bug-loop.asm}, \file{bug-stack.asm}, dan
          \file{program-larik.asm} yang disiapkan dosen.
\end{itemize}

\subsection{Langkah Kerja}

\begin{enumerate}
    \item \textbf{Menyiapkan berkas uji.}
    \begin{enumerate}
        \item Salin berkas \file{bug-loop.asm} ke direktori kerja
              DOSBox, lalu rakit dan gabungkan dengan perintah
              \file{tasm bug-loop.asm} dan \file{tlink bug-loop.obj}.
        \item Ulangi langkah yang sama untuk \file{bug-stack.asm}.
        \item Tuliskan program \eng{off-by-one} pada
              Kode~\ref{lst:off-by-one} ke dalam berkas
              \file{program-larik.asm}, lalu rakit.
    \end{enumerate}

    \item \textbf{Kasus 1: perulangan tak berujung
          (\file{bug-loop.exe}).}
    \begin{enumerate}
        \item Jalankan DEBUG, lalu muat program dengan urutan
              \code{N} dan \code{L}.
        \item Catat keluaran \code{R}. Berapa nilai \reg{CS}, \reg{IP},
              \reg{DS}, \reg{SS}, dan \reg{SP}?
        \item Tampilkan hasil pembongkaran dengan \code{U 0000 001c},
              lalu salin ke dalam laporan.
        \item Jalankan \code{T} sekurang-kurangnya sepuluh kali.
              Catat setiap langkah pada tabel \eng{trace}
              (gunakan susunan kolom pada
              Tabel~\ref{tab:trace-register}).
        \item Jawab: pada langkah berapa penyimpangan mulai terjadi,
              dan register apa yang menyimpang?
        \item Tuliskan hipotesis penyebab, buktikan dengan data
              \eng{trace}, lalu ubah program dan rakit ulang sampai
              keluarannya benar (\hex{000F} atau \dec{15}).
    \end{enumerate}

    \item \textbf{Kasus 2: kesalahan \eng{off-by-one}
          (\file{program-larik.exe}).}
    \begin{enumerate}
        \item Muat program ke dalam DEBUG, lalu jalankan sampai
              perulangan selesai menggunakan \code{G} pada alamat
              setelah \asm{LOOP}.
        \item Tampilkan isi larik dengan \code{D DS:0}.
        \item Tentukan berapa byte yang seharusnya dibaca dan berapa
              byte yang benar-benar dibaca.
        \item Perbaiki program, rakit ulang, dan buktikan bahwa
              hasilnya sekarang \hex{0096}.
    \end{enumerate}

    \item \textbf{Kasus 3: \eng{stack} tidak seimbang
          (\file{bug-stack.exe}).}
    \begin{enumerate}
        \item Muat program, lalu jalankan dengan \code{G 0005} agar
              berhenti tepat sebelum \asm{CALL}.
        \item Telusuri dengan \code{T} sampai seluruh register selesai
              ditekan. Perhatikan nilai \reg{SP} pada setiap langkah
              dan catat pada tabel \eng{trace} dengan menambahkan
              kolom \reg{SP}.
        \item Tampilkan \eng{stack} dengan \code{D SS:SP}, lalu
              gambarkan susunan bingkai \eng{stack}-nya seperti pada
              Gambar~\ref{fig:stack-frame}.
        \item Lanjutkan penelusuran sampai \asm{RET} dijalankan.
              Alamat berapa yang diambil oleh \asm{RET}, dan mengapa
              demikian?
        \item Bandingkan pemakaian \code{T} dan \code{P} pada
              instruksi \asm{CALL} tersebut. Tuliskan perbedaan
              keluarannya dan jelaskan kapan masing-masing lebih
              tepat dipakai.
        \item Perbaiki program sehingga pesan \code{Selesai.}
              tercetak.
    \end{enumerate}
\end{enumerate}

\subsection{Tabel Kerja}

Salin dan lengkapi tabel berikut pada jurnal praktikum Anda.

\begin{table}[H]
    \centering
    \caption{Templat tabel \eng{trace} untuk diisi mahasiswa.}
    \label{tab:template-trace}
    \scriptsize
    \setlength{\tabcolsep}{2.5pt}
    \begin{tabular}{@{}rl l rrrr r ccc p{2.3cm}@{}}
        \toprule
        \textbf{L} & \textbf{Alamat} & \textbf{Instruksi} &
        \textbf{AX} & \textbf{BX} & \textbf{CX} & \textbf{DX} &
        \textbf{SP} & \textbf{ZF} & \textbf{CF} & \textbf{OF} &
        \textbf{Keterangan} \\
        \midrule
        1 & & & & & & & & & & & \\
        \addlinespace
        2 & & & & & & & & & & & \\
        \addlinespace
        3 & & & & & & & & & & & \\
        \addlinespace
        4 & & & & & & & & & & & \\
        \addlinespace
        5 & & & & & & & & & & & \\
        \bottomrule
    \end{tabular}
\end{table}

\subsection{Hasil yang Dikumpulkan}

\begin{enumerate}
    \item Tabel \eng{trace} lengkap untuk ketiga kasus, minimal
          sepuluh langkah per kasus, memuat kolom \reg{AX},
          \reg{BX}, \reg{CX}, \reg{DX}, \reg{SP}, \flag{ZF},
          \flag{CF}, dan \flag{OF}.
    \item Salinan keluaran \code{R}, \code{U}, \code{D}, dan
          \code{T} untuk setiap kasus.
    \item Rumusan hipotesis dan bukti pengujiannya untuk setiap
          kasus, ditulis dengan alur metodologi pada
          Gambar~\ref{fig:metodologi-debugging}.
    \item Berkas hasil perbaikan ketiga program beserta tangkapan
          layar keluarannya.
    \item Perbandingan pengalaman memakai \code{T} dan \code{P} pada
          instruksi \asm{CALL}.
\end{enumerate}

\begin{asesmenbox}
\textbf{Asesmen kinerja \eng{debugging} (bobot 10\% dari komponen
praktikum).} Penilaian dilakukan terhadap laporan dan demonstrasi
langsung.

\begin{tabular}{@{}p{2.6cm}cp{9.6cm}@{}}
    \toprule
    \textbf{Aspek kinerja} & \textbf{Bobot} & \textbf{Kriteria} \\
    \midrule
    Penggunaan perintah DEBUG & 20\% &
    Mampu memuat berkas (\code{N}, \code{L}), menampilkan register
    (\code{R}), membongkar kode (\code{U}), dan menelusuri
    (\code{T}, \code{P}, \code{G}) tanpa bantuan \\
    \addlinespace
    Ketepatan tabel \eng{trace} & 25\% &
    Nilai register dan \eng{flag} tercatat benar pada setiap langkah;
    perbedaan \code{T} dan \code{P} dicatat dengan tepat \\
    \addlinespace
    Analisis \eng{dump} memori & 20\% &
    Mampu menafsirkan isi larik dan bingkai \eng{stack}; menemukan
    byte atau alamat yang menyimpang \\
    \addlinespace
    Ketepatan diagnosis & 20\% &
    Hipotesis sesuai bukti; penyebab yang disebutkan benar-benar
    sumber kesalahan, bukan gejala \\
    \addlinespace
    Mutu perbaikan dan verifikasi & 15\% &
    Perbaikan mengubah satu hal pada satu waktu dan seluruh kasus uji
    diulang setelah perbaikan \\
    \bottomrule
\end{tabular}
\end{asesmenbox}

\section{Latihan}
\label{sec:latihan-bab15}

\begin{enumerate}
    \item \textbf{Array satu dimensi.} Jelaskan cara mendeklarasikan
          array word lima elemen. Tentukan offset elemen pertama,
          kedua, dan kelima.

    \item \textbf{Array multi dimensi.} Diketahui matriks word
          $3 \times 4$. Tentukan offset elemen baris 2 kolom 1 dan
          baris 1 kolom 3.

    \item \textbf{String.} Jelaskan representasi string di memori.
          Mengapa string untuk layanan DOS fungsi \hex{09} harus
          diakhiri \code{'\$'}?

    \item \textbf{Menyalin string.} Tuliskan perulangan yang
          menyalin string dari \reg{SI} ke \reg{DI} sampai menemukan
          \code{'\$'}.

    \item \textbf{Membandingkan string.} Tuliskan perulangan yang
          membandingkan dua string karakter demi karakter. Kapan
          perbandingan dapat dihentikan?

    \item \textbf{Desimal.} Jelaskan algoritme mengubah bilangan
          menjadi digit desimal. Mengapa sisa pembagian disimpan
          pada \eng{stack}?

    \item \textbf{Heksadesimal.} Jelaskan algoritme mengubah
          bilangan menjadi empat digit heksadesimal. Mengapa nilai
          10--15 perlu ditambah 7 sebelum diubah menjadi karakter?

    \item \textbf{Penulisan program.} Tuliskan program yang
          menjumlahkan array word lima elemen, lalu mencetak hasilnya
          sebagai bilangan desimal. Gunakan \asm{LOOP} dan prosedur
          konversi desimal.
\end{enumerate}

\begin{refleksibox}
\vspace{1.1em}
Jawab pertanyaan berikut pada jurnal praktikum Anda.
\begin{enumerate}
    \item Bagian mana yang paling sulit: perhitungan offset array,
          penyalinan string, atau konversi bilangan? Mengapa?
    \item Setelah menelusuri konversi desimal, apakah Anda memahami
          mengapa \eng{stack} dipakai untuk membalik urutan digit?
          Jelaskan.
    \item Konsep apa yang perlu Anda ulangi sebelum mempelajari
          debugging level instruksi pada bab berikutnya?
\end{enumerate}
\end{refleksibox}

\section{Rangkuman}
\label{sec:rangkuman-bab15}

Pada bab ini telah dibahas:

\begin{itemize}
    \item Array adalah kumpulan data bertipe sama yang disimpan
          berurutan di memori. Array satu dimensi diakses dengan
          register penunjuk dan perulangan.

    \item Array multi dimensi disimpan baris demi baris. Offset
          elemen dihitung dengan
          $(r \times \text{KOLOM} + c) \times \text{ukuran elemen}$.

    \item String adalah array byte yang berisi kode ASCII. Akhir
          string dapat ditandai \code{'\$'} atau dengan menyimpan
          panjangnya.

    \item Penyalinan string dilakukan byte per byte sampai penanda
          akhir. Perbandingan dan pencarian memakai \asm{CMP} di
          dalam perulangan.

    \item Bilangan desimal diubah menjadi digit dengan pembagian
          berulang oleh 10. Sisa pembagian disimpan pada
          \eng{stack} agar urutan cetaknya benar.

    \item Bilangan heksadesimal diubah dengan mengambil empat bit
          sekaligus melalui \asm{ROL} dan \asm{AND} \hex{0F}.

    \item Digit 0--9 diubah menjadi karakter dengan menambahkan
          \hex{30}. Nilai 10--15 diubah menjadi \code{'A'}--\code{'F'}
          dengan menambahkan \hex{37}.

    \item Pencetakan dilakukan dengan layanan DOS
          \code{INT 21h} fungsi \hex{02} untuk satu karakter.
\end{itemize}

\begin{checklistbox}
\vspace{1.1em}
Tandai setiap butir yang sudah Anda kuasai.

\begin{itemize}
    \item[\checkbox] Saya dapat mendeklarasikan array satu dimensi.
    \item[\checkbox] Saya dapat menghitung offset array multi
          dimensi.
    \item[\checkbox] Saya dapat mengakses array dengan \reg{SI} dan
          \asm{LOOP}.
    \item[\checkbox] Saya dapat menyalin, membandingkan, dan mencari
          string.
    \item[\checkbox] Saya dapat mengubah bilangan menjadi digit
          desimal.
    \item[\checkbox] Saya dapat mengubah bilangan menjadi digit
          heksadesimal.
    \item[\checkbox] Saya dapat mencetak hasil konversi dengan
          layanan DOS.
\end{itemize}

\vspace{0.4em}
Bila masih ada butir yang belum bertanda, ulangi bagian yang
bersesuaian sebelum melanjutkan ke bab berikutnya.
\end{checklistbox}

\begin{infobox}
\vspace{1.1em}
\textbf{Rujukan bab.} Materi pada bab ini diringkas dari
\cite{gunawan2019}, Bagian 10, hlm.~93--100, Bagian 14.1,
hlm.~127--131, dan Bagian 16, hlm.~154--158, serta \cite{sto2001},
BAB IX, hlm.~41--46, dan BAB XIX--XX, hlm.~110--121.
\end{infobox}

\section{Evaluasi}
\label{sec:evaluasi-bab15}

\subsection{Soal Pilihan Ganda}

\begin{enumerate}
    \item Perintah DEBUG yang menampilkan seluruh register dan
          \eng{flag} adalah \dots
          \begin{enumerate}
              \item \code{U} \item \code{D}
              \item \code{R} \item \code{T}
          \end{enumerate}

    \item Pada saat \lstasm|CALL HITUNG| ditelusuri, perintah yang
          \emph{tidak} masuk ke dalam prosedur \asm{HITUNG} adalah
          \dots
          \begin{enumerate}
              \item \code{T} \item \code{P}
              \item \code{U} \item \code{R}
          \end{enumerate}

    \item Keluaran baris \eng{flag} DEBUG yang menunjukkan bahwa
          \flag{ZF} bernilai satu dan \flag{CF} bernilai nol adalah
          \dots
          \begin{enumerate}
              \item \code{NV UP EI PL NZ NA PO NC}
              \item \code{NV UP EI PL ZR NA PE NC}
              \item \code{OV DN DI NG NZ AC PO CY}
              \item \code{NV UP EI NG NZ NA PO CY}
          \end{enumerate}

    \item Sebuah prosedur menekan \reg{AX}, \reg{BX}, dan \reg{CX},
          tetapi hanya memulihkan \reg{CX} dan \reg{BX}. Akibat yang
          paling mungkin terjadi adalah \dots
          \begin{enumerate}
              \item assembler menolak program karena jumlah
                    \asm{PUSH} dan \asm{POP} tidak seimbang
              \item \asm{RET} melompat ke alamat yang bukan alamat
                    kembali
              \item register \reg{AX} otomatis dipulihkan oleh DOS
              \item program berjalan benar karena \reg{AX} tidak
                    dipakai lagi
          \end{enumerate}

    \item Program selesai tanpa pesan kesalahan dari assembler dan
          sistem operasi, tetapi hasil perhitungannya salah. Jenis
          kesalahan ini adalah \dots
          \begin{enumerate}
              \item kesalahan sintaks \item kesalahan logika
              \item kesalahan segmentasi \item \eng{overflow}
                    \eng{stack}
          \end{enumerate}
\end{enumerate}

\subsection{Soal Uraian}

\begin{enumerate}
    \item Jelaskan perbedaan perintah \code{T} dan \code{P} ketika
          dihadapkan pada instruksi \asm{CALL} dan \code{INT 21h}.
          Sertakan satu situasi ketika \code{T} lebih tepat dan satu
          situasi ketika \code{P} lebih tepat.

    \item Diberikan keluaran \code{D SS:00F4} berikut.

\begin{session}[caption={Keluaran \code{D} pada \eng{stack} untuk
soal uraian 2.}]
13A7:00F4  00 00 00 00 00 00 C3 1A-08 00 00 00 00 00 00 00  ................
\end{session}

          \begin{enumerate}
              \item Tentukan isi \reg{SS}:\hex{00FA} dan
                    \reg{SS}:\hex{00FC}.
              \item Bila \reg{SP} bernilai \hex{00FA} dan instruksi
                    berikutnya adalah \asm{RET}, alamat berapa yang
                    akan dimuat ke \reg{IP}? Jelaskan perhitungannya.
              \item Sebutkan satu penyebab yang paling mungkin bila
                    nilai tersebut bukan alamat kembali yang benar.
          \end{enumerate}

    \item Uraikan enam tahap metodologi \eng{debugging} sistematis.
          Untuk setiap tahap, sebutkan bukti yang harus dikumpulkan
          dan jelaskan mengapa tahap tersebut tidak boleh dilewati.

    \item Sebuah program menjumlahkan delapan byte larik, tetapi
          hasilnya selalu lebih besar dari yang seharusnya. Tuliskan
          langkah \eng{debugging} yang Anda lakukan mulai dari
          reproduksi sampai verifikasi. Sebutkan perintah DEBUG yang
          dipakai pada setiap langkah dan bukti yang diharapkan dari
          keluaran perintah tersebut.
\end{enumerate}

\subsection{Kisi-kisi Penilaian}

\begin{table}[H]
    \centering
    \caption{Kisi-kisi evaluasi Bab~\ref{chap:debugging}.}
    \label{tab:kisi-bab15}
    \small
    \begin{tabular}{@{}p{1.6cm}p{4.8cm}p{3.0cm}c@{}}
        \toprule
        \textbf{Sub-CPMK} & \textbf{Indikator yang diukur} &
        \textbf{Bentuk soal} & \textbf{Bobot} \\
        \midrule
        4.1 & Mengenali perintah penelusuran dan pemantauan \eng{flag} &
        Pilihan ganda 1, 3 & 15\% \\
        \addlinespace
        4.1 & Membedakan \code{T} dan \code{P} pada \asm{CALL} dan
        \code{INT} & Pilihan ganda 2; Uraian 1 & 25\% \\
        \addlinespace
        4.2 & Menganalisis \eng{dump} memori dan bingkai \eng{stack} &
        Uraian 2 & 25\% \\
        \addlinespace
        4.2 & Menjelaskan metodologi \eng{debugging} dan menerapkannya &
        Pilihan ganda 4, 5; Uraian 3, 4 & 35\% \\
        \bottomrule
    \end{tabular}
\end{table}

\begin{checklistbox}
\textbf{Checklist Pencapaian Kompetensi Bab~\ref{chap:debugging}}

Seorang mahasiswa dinyatakan kompeten pada Sub-CPMK 4.1 dan 4.2
apabila:
\begin{itemize}
    \item[\checked] dapat menyebutkan taksonomi kesalahan program
          beserta gejala khas masing-masing;
    \item[\checked] dapat memakai perintah \code{R}, \code{U},
          \code{T}, \code{P}, \code{G}, dan \code{D} tanpa panduan;
    \item[\checked] dapat menjelaskan kapan \code{T} dan kapan
          \code{P} dipakai pada \asm{CALL} serta \code{INT};
    \item[\checked] dapat mengisi tabel \eng{trace} register dan
          \eng{flag} dengan benar sekurang-kurangnya sepuluh langkah;
    \item[\checked] dapat membaca \eng{dump} memori dan menentukan
          isi bingkai \eng{stack};
    \item[\checked] dapat menemukan dan memperbaiki perulangan tak
          berujung, kesalahan \eng{off-by-one}, dan \eng{stack} tidak
          seimbang dalam batas waktu yang diberikan;
    \item[\checked] dapat menerapkan enam tahap metodologi
          \eng{debugging} dan menuliskan laporan \eng{trace} yang
          dapat diulang oleh orang lain.
\end{itemize}
\end{checklistbox}


\end{document}
